\subsection{A comparison from the source at each time}\label{sec:transport}
The same centered source estimate that locates the zeros in Section~\ref{sec:anchors} gives a comparison for their motion throughout a high-time interval. This calculation uses the phase velocity of Section~\ref{sec:current}.

\begin{theorem}[Signed backward transfer]\label{thm:transfer}
Fix \(C>0\). There is a time \(t_C\) such that, for \(T\ge t\ge t_C\) and a marked edge satisfying \(q_{k+1}\le C\inf_{t\le r\le T}\ell_r\), the same two positive phase levels are simple throughout. Uniformly on that family,
\begin{align}
 (\log g)^{\displaystyle\cdot}&=-2/\kappa_r+O_C(\kappa_r^{-2})<0,
 &M&=\frac1{2D}+\frac1{\kappa_r}+O_C(\kappa_r^{-2}),\label{eq:asymM}\\
 D(t)&=D(T)(L_T/L_t)^2
           \exp\!\left(O_C\!\left(\int_t^T\kappa_r^{-2}\dd r\right)\right).
 \label{eq:asymD}
\end{align}
For any \(X_r\) containing the endpoints in \([0,X_r]\) and satisfying \(X_r\le C'\sqrt{\kappa_r}\), there is also a mean-normalized bound
\begin{equation}
 \left|\log\frac{a_t g(t)}{a_Tg(T)}\right|
 \le C''\int_t^T X_r^2\kappa_r^{-3}\dd r. \label{eq:meanD}
\end{equation}
The constants depend only on the fixed window. The asymptotic theorem gives existence of \(t_C\); Section~\ref{sec:finite} treats a specific interval independently.
\end{theorem}
\begin{proof}
At a physical spatial argument \(x=\sqrt\kappa\,\zeta\), the analytic continuation of the normalized current \eqref{eq:current} is
\[
 j(t,\sqrt\kappa\,\zeta)
 =\exp\{K_0(\ell\zeta)+K_0(-\ell\zeta)\}\,[1-\Psi'(\ell\zeta)].
\]
This follows on the real axis from the amplitude and phase in \eqref{eq:odd} and then by continuation. The odd cubic in \eqref{eq:complex-saddle} cancels from the exponential. Cauchy's estimate on a slightly larger disk gives \(\Psi'(\ell\zeta)=O_C(1/(\ell\sqrt\kappa))\). Hence on every fixed complex disk, for large enough \(t\),
\begin{equation}
 \log j(t,\sqrt\kappa\,\zeta)=-\zeta^2+O_C(\kappa^{-1}).\label{eq:logj}
\end{equation}
The branch is real at zero. Two spatial derivatives give
\((\log j)_{xx}=-2/\kappa+O_C(\kappa^{-2})\). The inverse theorem places the full real edge in a fixed multiple of \([0,\sqrt\kappa]\) and identifies its index at each time; uniqueness and simplicity give the same moving endpoints. Average the curvature in \eqref{eq:velocity} and use \eqref{eq:forcing}. From \eqref{eq:saddle}, \(\dot L=2L/\kappa\). Integrating the estimate for \(2\mathcal A\) gives \eqref{eq:asymD}. The integrated remainder is \(O_C(1/[t(\log t)^2])\) even if the upper time tends to infinity.

Four spatial derivatives of \eqref{eq:logj} give \((\log j)_{xxxx}=O_C(\kappa^{-3})\). Evenness and \eqref{eq:difference} give \((\log j)_{xxx}(0)=0\) and \((\log j)_{xx}(0)=-\dot a/a\). Two real integrations therefore give \(|(\log j)_{xx}+\dot a/a|\ll_C x^2/\kappa^3\). Average over the edge and integrate in time, proving \eqref{eq:meanD}. For fixed \(k\), \(X_r=O_k(1/a_r)\) is allowed and the last integral is \(O_k(1/[t^2(\log t)^5])\).
\end{proof}
These are estimates on the complete interaction and the signed drift of the actual edge. The individual source expansions also localize its endpoints throughout the same high-time domain. The estimate establishes direction of motion and its integrated size there. Extending it to an earlier region requires additional source information on the intervening path.

\paragraph{The correction after removing the mean spacing.}
Let \(h=\log j\) and denote spatial averaging over the marked interval by
angle brackets. Combining \eqref{eq:velocity} with \(h_{xx}(t,0)=-\dot a/a\)
and \(h_{xxx}(t,0)=0\) gives the exact identity
\begin{equation}
 \frac{\dd}{\dd t}\log(ag)=
  \left\langle\int_0^x(x-y)\,\partial_y^4h(t,y)\dd y\right\rangle,
 \qquad \partial_x^4h(t,0)=m_4^*-3w^2 .\label{eq:quartic}
\end{equation}
Freezing the fourth derivative at zero yields
\((m_4^*-3w^2)(x_-^2+x_-x_++x_+^2)/6\).
If \(|\partial_x^6 h|\le B_6\) on \([0,X]\), evenness gives an
absolute error at most \(B_6 X^4/24\). This retains the cancellation between
the fourth moment and the squared variance when one compares a change
of gap with the change of its natural mean spacing.
